\bookchapter{פרק 63: תצוגות, מצב וניווט}

\chapterintro{כבר יודעים לבנות צעצוע של עמוד אחד. הפרק הזה פותר שתי בעיות של “גדילה”: איך מדלגים בין מספר מסכים (ניווט), ואיך נתונים זורמים בין המסכים (ניהול מצב). אחרי הפרק הזה תוכלו להקים שלד קלאסי של אפליקציה מרובת מסכים: “רשימה פרטים עריכה”.}

\booksection{“התמונה המלאה” של זרימת הנתונים בניהול מצב}

ב-SwiftUI כלי ניהול המצב לא מסתכמים ב-@State. קודם תסתכלו על התמונה המלאה ותבינו “מתי משתמשים במה”:

\begin{center}
\begin{tabularx}{\linewidth}{|X|X|X|X|}
\hline
\textbf{כלי} & \textbf{על מה מפקח} & \textbf{מי משתמש בו} & \textbf{תרחיש טיפוסי} \tabularnewline
\hline
@State & מצב זמני פרטי למסך & התצוגה הנוכחית & מונה, טקסט בשדה קלט \tabularnewline
@Binding & השאלת מצב מאחר & תצוגת בת משנה את תצוגת האב & רכיב מותאם אישית \tabularnewline
@StateObject / @Observable & לוגיקה משותפת בין מסכים & מודל תצוגה (ViewModel) & טעינת נתונים, לוגיקה עסקית \tabularnewline
@EnvironmentObject / .environment & מצב משותף גלובלי & כל האפליקציה & מצב התחברות המשתמש, הגדרות עיצוב \tabularnewline
\hline
\end{tabularx}
\end{center}

\begin{center}\small\color{gray}תמונת מצב מלאה של כלי המצב ב-SwiftUI (2026; @Observable כבר הבחירה המובילה)\end{center}

\begin{lstlisting}[language=Swift]
// הדרך המודרנית של iOS 17+: המאקרו @Observable (מומלץ לפרויקטים חדשים)
import Observation

@Observable
class CounterModel {                 // מודל התצוגה
    var count = 0
    func increment() { count += 1 }
}

struct CounterView: View {
    @State private var model = CounterModel()   // מחזיק את המודל

    var body: some View {
        VStack {
            Text("לחצת \(model.count) פעמים")
            Button("לחץ כאן") { model.increment() }
        }
    }
}
\end{lstlisting}

עצה בוטה למתחילים: למצב פרטי למסך — @State; למה שחוצה מסכים או כולל לוגיקה — מחלקת @Observable. במדריכים ישנים תמצאו @StateObject/@ObservableObject — זוהי הגישה של הדור הקודם; מספיק להבין אותה, ובקוד חדש ללכת לפי מה שלמעלה.

\booksection{NavigationStack: ה“מחסנית” של מעברי מסך}

מ-iOS 16 ואילך מנווטים בין מסכים עם NavigationStack — כמו ערימת צלחות: דוחפים אחת פנימה, ובחזרה שולפים אחת-אחת:

\begin{lstlisting}[language=Swift]
struct HomeView: View {
    var body: some View {
        NavigationStack {
            List {
                NavigationLink("הפנקס של דני") {
                    DetailView(name: "הפנקס של דני")   // מסך היעד
                }
                NavigationLink("הפנקס של נועה") {
                    DetailView(name: "הפנקס של נועה")
                }
            }
            .navigationTitle("רשימת הפנקסים")   // כותרת גדולה למעלה
        }
    }
}

struct DetailView: View {
    let name: String
    var body: some View {
        Text("פרטי \(name)")
            .navigationTitle(name)
            .navigationBarTitleDisplayMode(.inline)
    }
}
\end{lstlisting}

NavigationLink הוא “קישור שאפשר ללחוץ עליו ולעבור מסך”; navigationTitle יוצר אוטומטית כותרת לסרגל הניווט וכפתור חזרה. למעבר מתוכנת (למשל חזרה אוטומטית אחרי לחיצה על “שמור”) משתמשים ב-@Environment(\textbackslash{}.dismiss) לסגירת המסך הנוכחי — בשילוב שני הטריקים האלה מכוסים 90\% מצרכי הניווט.

\booksection{חלוניות קופצות ותגיות תחתונות: sheet ו-TabView}

מלבד “דחיפת מסך” יש שני אינטראקציות קלאסיות נוספות:

\begin{lstlisting}[language=Swift]
// 1. sheet שעולה מלמטה, חצי מסך או מסך מלא (האהובה על טפסי הזנה)
struct HomeView: View {
    @State private var showAdd = false
    var body: some View {
        Button("לרשום הוצאה") { showAdd = true }
            .sheet(isPresented: $showAdd) {
                AddRecordView()     // מסך ההזנה שנפתח
            }
    }
}

// 2. סרגל טאבים תחתון ("השלד" של האפליקציה)
struct MainTabView: View {
    var body: some View {
        TabView {
            RecordListView()
                .tabItem { Label("פנקס", systemImage: "book") }
            StatsView()
                .tabItem { Label("סטטיסטיקה", systemImage: "chart.pie") }
            SettingsView()
                .tabItem { Label("הגדרות", systemImage: "gear") }
        }
    }
}
\end{lstlisting}

\begin{tcbinfo}{טיפ: מתי משתמשים במה?}
כלל אצבע: את הזרימה הראשית דוחפים עם NavigationStack (רשימה פרטים), פעולות זמניות נפתחות בחלונית sheet (הוספה, עריכה, אישור), ובלוקים ראשיים שעומדים בפני עצמם — ב-TabView (תגיות תחתונות). השילוב של השלושה הוא השלד של רוב האפליקציות.

\end{tcbinfo}

\booksection{List: הרשימה “הלאומית” של iOS}

הרכיב הכי שכיח והכי “אפלי” ב-iOS הוא List — עם קווי הפרדה, מחיקה בהחלקה והדגשה בהקשה, מובנים מטבעם:

\begin{lstlisting}[language=Swift]
struct RecordListView: View {
    @State private var records = ["צהריים 25.5", "מטרו 6.0", "משכורת +8000"]

    var body: some View {
        List {
            ForEach(records, id: \.self) { r in
                Text(r)
            }
            .onDelete { indexSet in          // החלקה שמאלה למחיקה
                records.remove(atOffsets: indexSet)
            }
        }
    }
}
\end{lstlisting}

רשימה + NavigationLink + sheet + TabView — בשילוב ארבעת אלה תוכלו להקים כבר שלד אפליקציה שלם: “רשימת חשבונות פתיחת פרטים לחיצה על + לפתיחת חלונית הזנה”. פרויקט ספר החשבונות בפרק הבא נבנה בדיוק כך.

\begin{tcbpractice}{הידיים על ההגה: בונים "שלד של שלושה מסכים"}
\booksection{מחזור החיים של האפליקציה: חזית ורקע}

בדומה למחזור החיים של ה-Activity ב-Android, גם לאפליקציית iOS יש מעברי מצב של “לידה, חזית, רקע ומוות”. ב-SwiftUI עוקבים אחריהם בערך הסביבה scenePhase:

\begin{lstlisting}[language=Swift]
@main
struct MyApp: App {
    @Environment(\.scenePhase) private var phase

    var body: some Scene {
        WindowGroup {
            ContentView()
        }
        .onChange(of: phase) { _, newPhase in
            switch newPhase {
            case .active:    print("חזרה לחזית: לחדש את הרענון")
            case .inactive:  print("במעבר (שיחה או רגע המעבר לרקע)")
            case .background: print("מעבר לרקע: לשמור טיוטה")
            @unknown default: break
            }
        }
    }
}
\end{lstlisting}

שימושים טיפוסיים: רענון נתונים בחזרה לחזית, שמירת טיוטה בכניסה לרקע. iOS מפקח בקפדנות על ריצה ברקע — אפליקציה רגילה מושעית (מוקפאת) זמן קצר אחרי שיורדת לרקע, ואל תקוו לבצע עבודה כבדה בחשאי ברקע; למשימות מתוזמנות באמת השתמשו בהתראות המקומיות של פרק 68, ולא ב“ריצת רקע קבועה”.

בנו שלד של ספר חשבונות עם נתוני דמה: TabView עם שלוש תגיות (חשבונות/סטטיסטיקה/הגדרות) עמוד החשבונות מציג נתוני דמה ב-List, לחיצה על שורה עוברת למסך הפרטים (NavigationStack) כפתור “+” בפינה העליונה פותח חלונית הזנה ב-sheet. בלי שמירת נתונים ובלי אינטרנט — תרגול שלד טהור; כשהשלד הזה מוכן, מילוי הנתונים האמיתיים בפרק 65 יזרום מאליו.

\end{tcbpractice}

\begin{tcbsummary}{סיכום הפרק}
\begin{itemize}
  \item ארבעת כלי המצב: @State (פרטי למסך), @Binding (השאלה), @Observable (לוגיקה), .environment (גלובלי)
  \item בפרויקטים חדשים כותבים מודלי תצוגה עם המאקרו @Observable; את @ObservableObject מהמדריכים הישנים מספיק להבין
  \item NavigationStack דוחף מסכים; @Environment(\textbackslash{}.dismiss) לחזרה מתוכנת
  \item בזרימה הראשית דוחפים מסכים, פעולות זמניות נפתחות ב-sheet, ובלוקים ראשיים עוברים עם TabView
  \item List הוא "הרשימה הלאומית": קווי הפרדה ומחיקה בהחלקה מוכנים מהקופסה
\end{itemize}

\end{tcbsummary}
